\documentclass[11pt]{article}

% \providecommand{\answerDirectory}{solutions}

../../macros/macros.tex

% Local override, kept so this homework's PDF matches its original exactly:
% the shared macros (src/macros/macros.tex) now define answerBox differently (with an overflow check and a tick at the bottom).
\renewenvironment{answerBox}[1]{%
  \def\ANSWERBOXHEIGHT{#1}
  \begin{adjustbox}{valign=t}
  \begin{minipage}{\textwidth}
}{\end{minipage}\end{adjustbox}~~~\adjustbox{valign=t}{\parbox{0pt}{\color{gray}\rule[-\ANSWERBOXHEIGHT]{0.01pt}{\ANSWERBOXHEIGHT}}}}

\inputAnswer{problem_0_answer}

\doHwHeader{Homework \BLUE{9}}

%%%%% EDIT THE FILES NAMED problem_X_answer.tex.  DON'T EDIT THIS FILE

%%%%%%%%%%%%%%%%%%%%%%%%%%%%%%%%%%%%%%%%%%%%%%%%%%%%%%%%%%%% 
%%%%%%%%%%%%%%%%%%%%%%%%%%%%%%%%%%%%%%%%%%%%%%%%%%%%%%%%%%%% 
%%%%%%%%%%%%%%%%%%%%%%%%%%%%%%%%%%%%%%%%%%%%%%%%%%%%%%%%%%%% 
%%%%%%%%%%%%%%%%%%%%%%%%%%%%%%%%%%%%%%%%%%%%%%%%%%%%%%%%%%%% 

\begin{document}

\gradescopeInit


%%%%% EDIT THE FILES NAMED problem_X_answer.tex.  DON'T EDIT THIS FILE 

\paragraph{General instructions.}
% First read \BLUE{Lecture Note 5}, then
Answer the problems given in the rest of this document.
Edit this LaTeX template to add your answers, as described in the file named \texttt{00\_README.txt},
then compile it to produce \BLUE{\texttt{homework\_9.pdf}}\ifGradescope{} for submission to gradescope\fi.

\ifGradescope
For this homework, when you submit to gradescope you don't need to identify the pages
for each problem.  We're using gradescope's ``fixed-length'' grading mode.
To make it work, your answer for each part should fit in the allotted space,
so that each subsequent answer is on the expected page in the expected location.
Your PDF should have
\BLUE{9}
pages.  
\fi

\paragraph{Don't know?  Say so.}
If you don't know the answer to a problem or any part of the problem,
we encourage you to simply write \emph{``I don't know''} instead of giving an answer.
(If you write ``I don't know'' you can also add any questions you might have.)

\paragraph{Gentle reminder: homework is for learning.}
As noted in Lecture Note 1, engaging fully with the homeworks is the best way to learn the material.
\emph{The primary role of homeworks is to help you learn the material,
  and the primary role of feedback on homeworks is to help you improve your ideas.}

Please try to explain your ideas as clearly as you can.
The clearer your explanations are, the more useful the feedback on them can be.

\paragraph{Collaboration and citation policy.}
If you are doing this homework as part of a course, follow the course's collaboration and citation policy.
In any case, at the end of each answer, list the people you collaborated with,
and cite every source that your answer uses ideas from
(other than the lecture notes and the sources listed in the lecture notes).
If there are none, write ``none''.

%%%%% EDIT THE FILES NAMED problem_X_answer.tex.  DON'T EDIT THIS FILE 

\vfill
\noindent{\footnotesize\copyrightNotice}

%%% Local Variables:
%%% mode: latex
%%% TeX-master: "homework_2"
%%% End:


%%%%%%%%%%%%%%%%%%%%%%%%%%%%%%%%%%%%%%%%%%%%%%%%%%%%%%%%%%%% 
%%%%%%%%%%%%%%%%%%%%%%%%%%%%%%%%%%%%%%%%%%%%%%%%%%%%%%%%%%%% 
%%%%%%%%%%%%%%%%%%%%%%%%%%%%%%%%%%%%%%%%%%%%%%%%%%%%%%%%%%%% 
%%%%%%%%%%%%%%%%%%%%%%%%%%%%%%%%%%%%%%%%%%%%%%%%%%%%%%%%%%%% 

%%%%% EDIT THE FILES NAMED problem_X_answer.tex.  DON'T EDIT THIS FILE 

%%%%%%%%%%%%%%%% PROBLEM 1  %%%%%%%%%%%%%%%%%

\newpage 

\problem\emph{(Sections 7.5 and 7.6; Lecture Note 2)}

Consider the following problem, \prob{Counting Subset Sums}:

\begin{mylist}
\item[\bf input:] A pair $I=(v, T)$ where $v=(v_1, v_2, \ldots, v_n)$ is a sequence
  of positive integers and $T\ge 0$ is an integer.

\item[\bf output:]
  The number of subsequences of $v$ with total $T$.  That is,
  \[\Big|\big\{ S\subseteq\{1,2,\ldots,n\} : v(S) = T\big\}\Big|,\]
  where for notational convenience throughout we define
  $v(S) = \sum_{i\in S} v_i$.
\end{mylist}

This problem differs from \prob{Counting Ways to Make Change} (HW 7 Problem 3).
In this problem each of the given numbers can be used at most once.

\bigskip

Fix an input $(v, T)$.   For each $i\in\{0, 1,\ldots, n\}$ and $t\in\{0,1,\ldots,T\}$,
define $N(i, t)$ to be the number of subsequences of $(v_1, \ldots, v_i)$ with total $t$.
That is, 
\[N(i, t) = \Big|\big\{ S\subseteq\{1,2,\ldots,i\} : v(S) = t\big\}\Big|.\]

The correct output for the given input is $N(n, T)$.

\begin{enumerate}[label=(\alph*)]

\item
  Give a recurrence relation for $N(i, t)$ for all $i\in\{0,1,\ldots,n\}$ and $t\in\{0,1,\ldots, T\}$,
  including the base cases.
  State the recurrence relation as a lemma (see the template).

\item
  Complete the partial long-form proof of the lemma given in the template.
  You can use a proof structure similar to the one for the long-form proof of correctness for \prob{Knapsack} in Section 7.6.

\end{enumerate}

You don't need to give the resulting algorithm or analyze the running time.
But you can test your recurrence, by implementing the resulting algorithm,
at \url{https://www.hackerrank.com/nealeyoung-algs101-subset-sums-and-piles}
(the ``Counting Subset Sums'' challenge, which asks for the count mod $2^{30}$).
(This is just for testing. Hacker-Rank submissions won't be graded.)

% \newpage

\vfill

% \setlength{\ITEMHT}{1.9in}
% {% \injectEnumerate
\inputAnswer{problem_1_answer}
% }

\newpage

%%%%%%%%%%%%%%%%%%%%%%%%%%%%%%%%%%%%%%%%%%%%%%%%%%%%%%%%%%%% 
%%%%%%%%%%%%%%%%%%%%%%%%%%%%%%%%%%%%%%%%%%%%%%%%%%%%%%%%%%%% 
%%%%%%%%%%%%%%%%%%%%%%%%%%%%%%%%%%%%%%%%%%%%%%%%%%%%%%%%%%%% 
%%%%%%%%%%%%%%%%%%%%%%%%%%%%%%%%%%%%%%%%%%%%%%%%%%%%%%%%%%%% 

%%%%% EDIT THE FILES NAMED problem_X_answer.tex.  DON'T EDIT THIS FILE 

%%%%%%%%%%%%%%%% PROBLEM 2 %%%%%%%%%%%%%%%%%

\problem\emph{(Sections 7.5 and 7.6)} \GradescopeComment{This should be on page 4 (top) of your PDF.}

Design a dynamic-programming algorithm for \prob{Smallest Subset Sum}:
\begin{mylist}
\item[\bf input:] A pair $I=(v, T)$ where $v=(v_1,v_2, \ldots, v_n)$ is a sequence
  of positive integers and $T\ge 0$ is an integer.
\item[\bf output:]
  The minimum size of any subsequence of $v$ with total $T$:
  \[\min \big\{ |S| :\, S\subseteq\{1,2,\ldots,n\},\, v(S) = T\big\},\]
  where, for notational convenience throughout we define $v(S) = \sum_{i\in S} v_i$.
\end{mylist}
Recall that the minimum of the empty set is $\infty$.
 (So, if no subsequence of $v$ has total $T$, the output should be $\infty$.)

\begin{enumerate}[label=(\alph*)]

\item 
  Per Section 7.4,
  define the subproblems that your algorithm solves for a given input $(v=(v_1,\ldots, v_n), T)$.

\item Which subproblem gives the final answer?

\item State an appropriate recurrence relation, including boundary cases. 

\item In terms of $n$ and $T$, what is the big-$\Theta$ running time of the resulting iterative (or recursive and memoized) algorithm?

\end{enumerate}

You can test your algorithm at \url{https://www.hackerrank.com/nealeyoung-algs101-subset-sum-and-lcs}.
(This is just for testing. Hacker-Rank submissions won't be graded.)

\vfill

\setlength{\ITEMHT}{2in}
{\injectEnumerate
  \inputAnswer{problem_2_answer}
}


%%%%%%%%%%%%%%%%%%%%%%%%%%%%%%%%%%%%%%%%%%%%%%%%%%%%%%%%%%%% 
%%%%%%%%%%%%%%%%%%%%%%%%%%%%%%%%%%%%%%%%%%%%%%%%%%%%%%%%%%%% 
%%%%%%%%%%%%%%%%%%%%%%%%%%%%%%%%%%%%%%%%%%%%%%%%%%%%%%%%%%%% 
%%%%%%%%%%%%%%%%%%%%%%%%%%%%%%%%%%%%%%%%%%%%%%%%%%%%%%%%%%%% 

%%%%% EDIT THE FILES NAMED problem_X_answer.tex.  DON'T EDIT THIS FILE 

%%%%%%%%%%%%%%%% PROBLEM 3 %%%%%%%%%%%%%%%%%

\newpage

\problem\emph{(Section 7.8)}

Design an $O(\ell m n)$-time dynamic-programming algorithm for the \prob{Three-Way LCS} problem:
\begin{mylist}
\item[\bf input:] a triple $I=(A[1..\ell],B[1..m],C[1..n])$ of three sequences.
\item[\bf output:] the maximum length of any common subsequence of $A[1..\ell]$, of $B[1..m]$, and of $C[1..n]$.
  \smallskip 
  (That is, the maximum length of any sequence that is simultaneously a subsequence of $A$, of $B$, and of $C$.)
\end{mylist}

\begin{enumerate}[label=(\alph*)]
\item Per Section 7.4, define the subproblems the algorithm solves given input $(A[1..\ell], B[1..m], C[1..n])$.
\item Which subproblem gives the final answer? 
\item State the recurrence relation, including boundary cases.
\item What is the running time of the resulting (iterative, or recursive and memoized) algorithm?
\end{enumerate}

You can test your algorithm at \url{https://www.hackerrank.com/nealeyoung-algs101-subset-sum-and-lcs}.
(This is just for testing. Hacker-Rank submissions won't be graded.)

\vfill
\setlength{\ITEMHT}{2in}
{\injectEnumerate
  \inputAnswer{problem_3_answer}
}

%%%%%%%%%%%%%%%%%%%%%%%%%%%%%%%%%%%%%%%%%%%%%%%%%%%%%%%%%%%% 
%%%%%%%%%%%%%%%%%%%%%%%%%%%%%%%%%%%%%%%%%%%%%%%%%%%%%%%%%%%% 
%%%%%%%%%%%%%%%%%%%%%%%%%%%%%%%%%%%%%%%%%%%%%%%%%%%%%%%%%%%% 
%%%%%%%%%%%%%%%%%%%%%%%%%%%%%%%%%%%%%%%%%%%%%%%%%%%%%%%%%%%% 

%%%%% EDIT THE FILES NAMED problem_X_answer.tex.  DON'T EDIT THIS FILE 

%%%%%%%%%%%%%%%% PROBLEM 4 %%%%%%%%%%%%%%%%%

\newpage

\problem 

Consider the flow network $G=(V,E)$ with edge capacities as shown below:

\begin{center}
  \begin{tikzpicture}[
    vertex/.style={circle, inner sep=0pt, minimum width=15pt, draw}, 
    label/.style={midway, fill=white, inner sep=2pt}, 
    every path/.style={->}
    ]
    \node[vertex] (s) {s}; 
    \node[vertex] (b) [right =of s] {b}; 
    \node[vertex] (a) [above =of b] {a}; 
    \node[vertex] (c) [below =of b] {c}; 
    \node[vertex] (d) [right =of a, xshift=20pt] {d}; 
    \node[vertex] (f) [right =of b, xshift=20pt] {f}; 
    \node[vertex] (g) [right =of c, xshift=20pt] {g}; 
    \node[vertex] (t) [right =of f] {t}; 

    %% EDGES AND CAPACITIES START HERE

    \path (s) edge node [label] {10} (a); % edge s -> a
    \path (s) edge node [label] {5}  (b); % edge s -> b
    \path (s) edge node [label] {15}  (c); % etc.

    \path (a) edge node [label] {4} (b) ;
    \path (a) edge node [label] {9} (d) ;
    \path (a) edge node [label] {15} (f); 

    \path (b) edge node [label] {8} (f) ;
    \path (b) edge node [label] {4} (c) ; 

    \path (c) edge node [label] {30} (g) ; 

    \path (d) edge node [label] {10} (t) ; 
    \path (d) edge node [label] {15} (f) ; 

    \path (f) edge node [label] {10} (t) ; 
    \path (f) edge node [label] {15} (g) ; 

    \path (g) edge node [label] {10} (t) ; 
    \path (g) edge node [label] {6} (b) ; 
  \end{tikzpicture}
\end{center}

\begin{enumerate}[label=(\alph*)]

\item 
  Find a maximum-value $s$-$t$ flow $f$ in the network.
  Draw the network, labeling each edge $(u, w)$ with the flow $f(u,w)$ on the edge.

\item
  What is the value of the flow? 

\item
  Find an $s$-$t$ cut $(S, T=V\setminus S)$ of minimum capacity in the network.
  For the cut you found, what is $S$?

\item
  And what is the capacity of the cut?

\end{enumerate}

\vfill

\setlength{\ITEMHT}{1in}
{\injectEnumerate
  \inputAnswer{problem_4_answer}
}

\end{document}

%%% Local Variables:
%%% mode: latex
%%% TeX-master: t
%%% End:
